\begin{abstract}
Language models trained on second-language (L2) learner corpora develop
structured, acquisition-aligned error profiles rather than uniform
degradation of a native-English baseline. This suggests the models
internalise coarse interlanguage regularities, but it offers no explicit
control: at inference, the model cannot be asked to predict as an A1
Spanish-L1 learner rather than a C1 Japanese-L1 learner. We propose a
light architectural modification that adds such control by conditioning
the attention gate on learner metadata. Building on the Qwen team's
finding that a head-specific sigmoid gate applied at the scaled
dot-product attention output (\emph{SDPA-output gating}) is the most
effective single change to softmax attention, we extend the gate's input
from the hidden state alone to the concatenation of the hidden state
with learned CEFR-proficiency and L1 embeddings. This yields a
per-learner, per-token filter on attention outputs at every layer,
activated only when the learner profile changes. We describe a pilot
design on EFCAMDAT with zero-shot transfer to three external learner
corpora (andrew100k, CELVA-SP, KUPA-KEYS), and enumerate a set of
complementary experimental variants that ablate the gate site, the gate
granularity, the activation form, and the metadata encoding. We argue
this combination --- Qwen-optimal attention gating paired with
learner-metadata conditioning --- is a minimal-change, interpretable
handle on \emph{predict-the-missing-word} tasks in SLA.
\end{abstract}
